\begin{enumerate}[label=\textbf{\ansnum.\arabic*}]
  \item Anchor position: $\exp(-(0.1/0.3)^2)=0.895$; anchor orientation: $\exp(-(0.2/0.4)^2)=0.779$; body positions: $\exp(-(0.05/0.3)^2)=0.973$. Weighted sum: $0.5\cdot0.895+0.5\cdot0.779+1.0\cdot0.973=0.447+0.389+0.973=1.81$ (of $2.0$ possible for these three terms).
  \item The global position term of the anchor falls to $\exp(-(0.5/0.3)^2)=0.062$ (weight $0.5$). The terms for the relative body positions, the orientations and the velocities stay at (nearly) $1$, because they describe the \emph{shape} and the motion of the body with respect to the anchor, not where it is. The design separates ``where am I'' (global, a weak weight, so a small drift is tolerated) from ``how am I moving'' (relative, a large weight).
  \item If the feet or the torso are far from the reference, the episode is clearly failed. Without a termination the robot lies on the floor for the remaining $8$\,s and collects only useless, noisy experience (the reference keeps moving away); the rewards there are all near $0$ and give no information. A termination saves this time and also gives a signal: ``falling is bad'' (the lost future reward).
  \item $12\cdot50=600$ frames; $600/20=30$ frames per bin. $7.3$\,s is frame $365$, which is in bin $\lfloor365/30\rfloor=12$ (counting from $0$; the $13$th bin).
  \item $D=-1$: $0$; $D=0$: $0.75$; $D=0.5$: $0.9375$; $D=1$: $1$; $D=2$: $0.75$. The maximum is $1$ at $D=1$, the target value for reference data: an output \emph{above} the target is not better but just as far from it as an output below, so there is no reason for more reward. (The score is a bounded quadratic.)
  \item (a) $0.55\cdot0.6\cdot1+0.45\cdot0=0.33$. (b) $0.55\cdot0.6\cdot0+0.45\cdot1=0.45$. (c) $0.55\cdot0.6\cdot0.75+0.45\cdot0.5=0.2475+0.225=0.4725$. With these numbers the mixture prefers (b) to (a): reaching the goal in an unnatural way is worth more than a perfectly natural standing still. This is the intended trade; if you want a stronger style, raise the scale $c$ or lower $\lambda$.
  \item $23+23+1+3+3+3+12+12=80$, and the discriminator input is $160$. For the K1: $22+22+1+3+3+3+12+12=78$ (and $156$ for the two states).
  \item The style reward is $\max(0,1-0.25\cdot4)=0$ for every transition of the policy: the signal vanishes, and the policy has no direction to improve in the style term. A discriminator that is too sharp is a typical failure of adversarial training. The gradient penalty $10\,\mathbb E\|\nabla D\|^2$ on the reference data makes $D$ change slowly with its input, so that nearby transitions get nearby scores and there is a slope to follow; it also keeps the discriminator from becoming sharper than the data justify. (The learning rate and the number of discriminator updates are the other levers.)
  \item AMP does not need a time index: the discriminator only asks whether a pair of consecutive states is plausible, and the policy can switch between the styles in the data. Tracking has to know which frame of which clip is current, so it is one policy per clip. If the data contain only forward walking, a turn on the spot is not in the style distribution: the discriminator will score it low, and the policy has to choose between the task reward (turn) and the style reward (walk forward), according to the weights; typically the robot turns with small steps along an arc, or becomes unstable.
  \item A sketch (not run): the failure probability of bin $b$ becomes $p_b+0.5\,p_{b+1}$ (capped at $1$), the sampling probabilities are convolved with the kernel $(0.25,0.5,0.25)$ before normalising, and the practice $p\leftarrow p\,e^{-0.002n}$ is applied to the same bins. I expect that smoothing spreads the samples over the hard bins and their neighbours and that adaptive sampling is still much faster than uniform sampling; but I did not run this variant, so I do not report numbers.
\end{enumerate}
